\makeabstract
    {
    Dopaminergic Signaling Contributes to the Regulation of Territorial Competition in Betta splendens
    }
    {
    Nazario Ramos\textsuperscript{1}, Jahzara Wilson\textsuperscript{1}, Ruth Berganross\textsuperscript{1}, Elaine Aquino\textsuperscript{1}, Kelly Wallace\textsuperscript{1}
    }
    {
    \textsuperscript{1} Neuroscience Program, Amherst College, Amherst, MA
    }
    {
    This study determines whether male Betta splendens develop a conditioned place preference (CPP) for a mirror-associated positive stimulus environment and the role of dopamine D1 receptor signaling in their aggressive behavior and locomotion.
    }
    {
    A conditioned place preference (CPP) assay is a behavioral test used to evaluate the motivational and rewarding properties of a stimulus. In this study, male Betta splendens were exposed to either a mirror stimulus or a control condition for 10-15 minutes once daily over a 10-day conditioning period. Conditioned preference was assessed by measuring the amount of time fish spent on the stimulus-associated side of the testing tank. To investigate the role of dopamine signaling in aggression and associative learning, fish were treated with SCH, an antagonist drug that blocks dopamine activity, or SKF, an agonist drug that enhances dopamine activity. These tests were conducted to examine how changes in dopamine signaling influence Betta splendens aggressive behavior and locomotion. 
    }
    {
    The results of the conditioned place preference (CPP) test suggested that there was no statistically significant correlation in the betta fish choosing between their initial preferred side (IPS) and the {\textquotedblleft}reward{\textquotedblright} mirror stimulus. The data showed random learning habits, or lack thereof, amongst both control and mirror exposed fish. However, the results of the dopamine experiment indicated that there was a statistically significant difference between the latency to lateral display before and after the dopamine receptor activator (SKF) was administered to the fish. The data revealed that SKF significantly increased the latency after its application. These findings suggest that dopamine has some effect on competition and aggression. 
    }
    {
    The results from this study suggest that dopamine does in fact influence territorial competition by changing the Betta splendens{\textquoteright} perception of social reward. These findings provide insight into how the neural mechanisms underlying aggression and reward may contribute to a broader understanding of how dopamine may affect social behavior across vertebrates.  
    }

    \newpage
    